Human-centered coding agents are increasingly designed to ask clarifying questions
rather than guess at ambiguous requirements. A recent empirical study on
general-purpose agents~\citep{aspi2026} reports that adversarial content delivered
through a clarification channel amplifies attack success rates by 15--20$\times$
compared to standard tool-output injection. We ask whether this \emph{clarification
tax} replicates for coding agents, using a matched benchmark we introduce called
\textbf{ClarInject-Code}: 25 seed scenarios across five attack categories
(backdoor, secret exfiltration, safety-check removal, dependency poisoning,
destructive commands) with machine-checkable success predicates. Evaluating
Llama-3.1-8B-Instruct, we find a \emph{negative} clarification tax ($\tax = -0.32$):
the model is substantially more vulnerable through tool-output injection (ASR $=
0.76$) than through the clarification channel (ASR $= 0.44$), suggesting the ASPI
effect may be specific to frontier-scale models with strong instruction-hierarchy
training. We further find that a simple prompt-tagging defense (D1) \emph{increases}
clarification ASR by 17 points for 8B models---a tag-authority inversion effect
where XML delimiters signal trustworthiness rather than isolation. Structural defenses
(structured extraction, plan-diff gating) reduce clarification ASR to 0.08 and
0.013 respectively. We release ClarInject-Code and all evaluation code.
